\chapter{Intro_QM}
\label{chap:foundations_qm}

Quantum mechanics is the theoretical framework that describes the behaviour of microscopic systems such as electrons, atoms and photons. 
It emerged in the early twentieth century as a response to a number of experimental observations that could not be reconciled with classical physics: black-body radiation, the photoelectric effect, atomic spectra and the stability of matter, to name a few prominent examples. 
What eventually replaced the classical picture is a probabilistic theory in which states are represented by vectors in an abstract Hilbert space and physical quantities by operators acting on that space.

The goal of this chapter is not to provide a complete historical overview, but to introduce the mathematical structure and postulates of quantum theory in the language used throughout the rest of the thesis. 
We start from the concept of Hilbert space and pure states, then introduce mixed states and density operators, generalized measurements, and the most general form of physical evolutions, described by completely positive trace-preserving maps (quantum channels). 
Finally, we discuss in some detail the special case of qubits and spin-$1/2$ systems, which will play a central role in the following chapters.

Throughout this chapter we keep $\hbar$ explicit in the equations; in later parts of the thesis we will often set $\hbar = 1$ for convenience. 
We denote by $\mathbb{1}$ the identity operator on the relevant Hilbert space, and by $\mathbb{1}_A$, $\mathbb{1}_B$, etc.\ the identity on a specific subsystem.

\section{From classical to quantum}

\subsection{Classical states and observables}

In classical mechanics the state of a particle is specified by a point in phase space, labelled by canonical coordinates $(q,p)$. 
A classical observable is a real-valued function $f(q,p)$, and its value on a given state is simply $f$ evaluated at that point. 
Even when classical statistical ensembles are considered, probabilities describe our ignorance about the actual point in phase space, not a fundamental randomness of the theory.

Time evolution is deterministic and governed by Hamilton's equations,
\begin{equation}
	\frac{\mathrm{d}q}{\mathrm{d}t} = \frac{\partial H}{\partial p},
	\qquad
	\frac{\mathrm{d}p}{\mathrm{d}t} = -\frac{\partial H}{\partial q},
\end{equation}
or, equivalently, by the Liouville equation for the phase-space distribution. 
Conceptually, once the initial state is fixed, all future values of any observable are determined.

\subsection{Quantum principles and the need for a new framework}

Several experimental facts are incompatible with this picture: 
the ultraviolet catastrophe in black-body radiation, 
the discrete line spectra of atoms 
and the stability of matter all hint at an intrinsic discreteness and at the impossibility of assigning definite values to all physical quantities simultaneously.

The quantum theory that emerged incorporates a few key principles:
\begin{itemize}
	\item \emph{Superposition principle:} if $\ket{\psi_1}$ and $\ket{\psi_2}$ are possible states, then any linear combination $c_1\ket{\psi_1} + c_2\ket{\psi_2}$ is also a possible state.
	\item \emph{Incompatibility of observables:} certain pairs of observables, such as position and momentum, cannot be simultaneously sharp, leading to uncertainty relations.
	\item \emph{Intrinsic probabilistic character:} the theory assigns probabilities to measurement outcomes; these probabilities are not, in general, attributable to ignorance about an underlying classical state.
	\item \emph{Entanglement:} composite systems can exhibit correlations that do not admit a classical probabilistic explanation in terms of local hidden variables.
\end{itemize}
To accommodate these features, quantum mechanics describes states and observables within the abstract framework of Hilbert spaces and operators. 
We now move to this formalism and to a more systematic statement of the postulates of quantum theory.

\section{Hilbert spaces, states and observables}

\subsection{Hilbert space and pure states}

A quantum system is associated with a complex separable Hilbert space $\mathcal{H}$, i.e.\ a complex vector space endowed with an inner product $\braket{\cdot|\cdot}$ that is complete with respect to the associated norm. 
A \emph{pure state} of the system is represented by a unit vector $\ket{\psi} \in \mathcal{H}$, with
\begin{equation}
	\braket{\psi|\psi} = 1.
\end{equation}
Physical states are rays in Hilbert space, meaning that $\ket{\psi}$ and $e^{i\varphi}\ket{\psi}$ represent the same physical state for any real phase $\varphi$.

An orthonormal basis $\{\ket{n}\}$ of $\mathcal{H}$ satisfies
\begin{equation}
	\braket{m|n} = \delta_{mn},
	\qquad
	\sum_n \ket{n}\bra{n} = \mathbb{1}.
\end{equation}
Any state $\ket{\psi}$ can be expanded as
\begin{equation}
	\ket{\psi} = \sum_n c_n \ket{n},
	\qquad
	c_n = \braket{n|\psi},
\end{equation}
and the normalization condition reads $\sum_n |c_n|^2 = 1$.

\subsection{Observables and spectral decomposition}

A physical observable $X$ is represented by a self-adjoint operator $X = X^\dagger$ on $\mathcal{H}$. 
By the spectral theorem, $X$ admits a decomposition of the form
\begin{equation}
	X = \sum_x x \, P_x,
\end{equation}
where $x$ are the real eigenvalues of $X$ and $P_x$ are orthogonal projectors onto the corresponding eigenspaces,
\begin{equation}
	P_x P_{x'} = \delta_{xx'} P_x,
	\qquad
	\sum_x P_x = \mathbb{1}.
\end{equation}
In the case of a continuous spectrum, such as position or momentum, the sum is replaced by an integral and the projectors $P_x$ become projection-valued measures.

\subsection{Composite systems and tensor product structure}

For a composite system made of two subsystems $A$ and $B$, with Hilbert spaces $\mathcal{H}_A$ and $\mathcal{H}_B$, the total system is described by the tensor product space
\begin{equation}
	\mathcal{H}_{AB} = \mathcal{H}_A \otimes \mathcal{H}_B.
\end{equation}
If $\{\ket{a}\}$ and $\{\ket{b}\}$ are orthonormal bases of $\mathcal{H}_A$ and $\mathcal{H}_B$, then $\{\ket{a}\otimes\ket{b}\}$ is an orthonormal basis of $\mathcal{H}_{AB}$. 
For brevity we write $\ket{a}\otimes\ket{b} = \ket{ab}$.

Operators acting on only one subsystem are extended to the full space by tensoring with the identity on the other subsystem, for example
\begin{equation}
	X_A \mapsto X_A \otimes \mathbb{1}_B,
	\qquad
	Y_B \mapsto \mathbb{1}_A \otimes Y_B.
\end{equation}

\subsection{Postulates of quantum mechanics (formal version)}

The discussion above can be condensed into a standard list of postulates:

\begin{itemize}
	\item \textbf{Postulate 1 (states).}
	The state of a quantum system is represented by a unit vector $\ket{\psi}$ in a complex Hilbert space $\mathcal{H}$, or, more generally, by a density operator $\hat{\rho}$ (see Sec.~\ref{sec:density_mixed}).
	
	\item \textbf{Postulate 2 (observables).}
	Every physical observable is associated with a self-adjoint operator $X$ on $\mathcal{H}$. 
	The only possible outcomes of a measurement of $X$ are the eigenvalues of $X$.
	
	\item \textbf{Postulate 3 (measurement).}
	If the system is in the state $\hat{\rho}$ and a projective measurement of $X = \sum_x x P_x$ is performed, the probability of obtaining the outcome $x$ is $p(x) = \operatorname{Tr}[\hat{\rho} P_x]$. 
	Immediately after the measurement, conditioned on the outcome $x$, the state collapses to
	\begin{equation}
		\hat{\rho} \longrightarrow \frac{P_x \hat{\rho} P_x}{p(x)}.
	\end{equation}
	More general measurements are described by positive operator-valued measures (POVMs) (see Sec.~\ref{sec:povm}).
	
	\item \textbf{Postulate 4 (dynamics).}
	If the system is closed, the time evolution of the state is unitary and governed by the Schrödinger equation. 
	In the density operator formalism this becomes the Liouville--von Neumann equation
	\begin{equation}
		i \hbar \frac{\mathrm{d}\hat{\rho}}{\mathrm{d}t} = [H,\hat{\rho}],
	\end{equation}
	where $H$ is the Hamiltonian operator. 
	More general, non-unitary evolutions arise when the system interacts with an environment and are described by completely positive trace-preserving maps (see Sec.~\ref{sec:channels}).
	
	\item \textbf{Postulate 5 (composite systems).}
	The Hilbert space of a composite system is the tensor product of the Hilbert spaces of the subsystems. 
	States and observables of subsystems are obtained by partial trace and by tensoring with the identity operator, respectively.
\end{itemize}

\section{Mixed states and density operators}
\label{sec:density_mixed}

\subsection{Statistical mixtures and density operators}

The state vector description is complete for pure states, but in many situations we deal with classical uncertainty about the preparation of the system. 
Suppose the system is prepared in the pure state $\ket{\psi_k}$ with probability $p_k$. 
The ensemble $\{p_k, \ket{\psi_k}\}$ is described by the density operator
\begin{equation}
	\hat{\rho} = \sum_k p_k \, \ket{\psi_k}\bra{\psi_k}.
\end{equation}
By construction,
\begin{equation}
	\hat{\rho} \ge 0,
	\qquad
	\operatorname{Tr}[\hat{\rho}] = 1.
\end{equation}
Conversely, any positive semidefinite trace-one operator on $\mathcal{H}$ can be interpreted as the density operator of some ensemble of pure states.

The expectation value of an observable $X$ in the state $\hat{\rho}$ is given by
\begin{equation}
	\langle X \rangle = \operatorname{Tr}[\hat{\rho} X],
\end{equation}
and the probability of obtaining the outcome $x$ in a projective measurement of $X = \sum_x x P_x$ is
\begin{equation}
	p(x) = \operatorname{Tr}[\hat{\rho} P_x].
\end{equation}
Pure states correspond to one-dimensional projectors $\hat{\rho} = \ket{\psi}\bra{\psi}$ and satisfy $\operatorname{Tr}[\hat{\rho}^2] = 1$, while mixed states have
\begin{equation}
	\operatorname{Tr}[\hat{\rho}^2] < 1.
\end{equation}

\subsection{Reduced states and partial trace}

Consider again a bipartite system $AB$ with Hilbert space $\mathcal{H}_A \otimes \mathcal{H}_B$ and global density operator $\hat{\rho}_{AB}$. 
If we are interested only in the statistics of observables acting on subsystem $A$, the relevant state is the \emph{reduced density operator} on $A$, defined by the \emph{partial trace} over $B$:
\begin{equation}
	\hat{\rho}_A = \operatorname{Tr}_B[\hat{\rho}_{AB}].
\end{equation}
Formally, the partial trace is the unique operation satisfying
\begin{equation}
	\operatorname{Tr}_{AB}\bigl[\hat{\rho}_{AB}(X_A \otimes \mathbb{1}_B)\bigr]
	=
	\operatorname{Tr}_A[\hat{\rho}_A X_A]
\end{equation}
for all observables $X_A$ on $\mathcal{H}_A$. 
An analogous definition yields $\hat{\rho}_B = \operatorname{Tr}_A[\hat{\rho}_{AB}]$.

An important consequence is that even if the global state $\hat{\rho}_{AB}$ is pure, the reduced state of a subsystem can be mixed. 
For instance, for the Bell state
\begin{equation}
	\ket{\Phi^+} = \frac{1}{\sqrt{2}}\bigl(\ket{00} + \ket{11}\bigr),
\end{equation}
we have $\hat{\rho}_{AB} = \ket{\Phi^+}\bra{\Phi^+}$, but
\begin{equation}
	\hat{\rho}_A = \hat{\rho}_B = \frac{\mathbb{1}_2}{2},
\end{equation}
which is maximally mixed.

\subsection{Purification and environment}

Every mixed state can be seen as arising from a pure state on a larger Hilbert space. 
Given $\hat{\rho}$ on $\mathcal{H}_S$, there exists an auxiliary Hilbert space $\mathcal{H}_E$ and a pure state $\ket{\Psi}_{SE} \in \mathcal{H}_S \otimes \mathcal{H}_E$ such that
\begin{equation}
	\hat{\rho} = \operatorname{Tr}_E \bigl[\ket{\Psi}_{SE}\bra{\Psi}_{SE}\bigr].
\end{equation}
The state $\ket{\Psi}_{SE}$ is called a \emph{purification} of $\hat{\rho}$. 
Purifications are not unique: if $\ket{\Psi_{SE}}$ and $\ket{\Phi_{SE}}$ are purifications of the same density operator, they are related by a unitary operator acting only on the ancilla,
\begin{equation}
	\ket{\Phi}_{SE} = (\mathbb{1}_S \otimes U_E) \ket{\Psi}_{SE}.
\end{equation}
This viewpoint is particularly useful to describe open systems: the environment can be modelled as an additional quantum system, and the apparent mixedness of the system's state is interpreted as arising from entanglement with the environment.

\subsection{Purity and von Neumann entropy}

A simple quantifier of mixedness is the \emph{purity},
\begin{equation}
	\gamma(\hat{\rho}) = \operatorname{Tr}[\hat{\rho}^2].
\end{equation}
We have $0 < \gamma(\hat{\rho}) \leq 1$, with $\gamma = 1$ if and only if $\hat{\rho}$ is pure. 
Another important measure is the von Neumann entropy,
\begin{equation}
	S(\hat{\rho}) = - \operatorname{Tr}[\hat{\rho} \log \hat{\rho}],
\end{equation}
which generalizes the Shannon entropy to the quantum domain. 
Pure states have $S(\hat{\rho}) = 0$, while mixed states have $S(\hat{\rho}) > 0$.

In the context of open systems, loss of purity and increase of entropy are associated with decoherence and with the build-up of correlations with the environment.

\section{Quantum measurements}
\label{sec:measurements}

\subsection{Projective measurements and the Born rule}

In the simplest measurement model, an ideal measurement of an observable $X = \sum_x x P_x$ is associated with the set of projectors $\{P_x\}$. 
If the system is prepared in the state $\hat{\rho}$, the probability of obtaining the outcome $x$ is
\begin{equation}
	p(x) = \operatorname{Tr}[\hat{\rho} P_x],
\end{equation}
and, conditioned on observing $x$, the post-measurement state is
\begin{equation}
	\hat{\rho} \longrightarrow \hat{\rho}_x =
	\frac{P_x \hat{\rho} P_x}{p(x)}.
\end{equation}
The unconditional post-measurement state, when the outcome is ignored, is obtained by averaging over all outcomes:
\begin{equation}
	\hat{\rho}' = \sum_x P_x \hat{\rho} P_x.
\end{equation}

\subsection{Incompatible observables and uncertainty relations}

If two observables $X$ and $Y$ commute, $[X,Y] = 0$, they admit a common eigenbasis and can be measured simultaneously with arbitrary precision. 
If they do not commute, they are called \emph{incompatible} and are subject to uncertainty relations.

For a state $\hat{\rho}$, the variance of an observable $X$ is 
\begin{equation}
	\Delta^2 X = \langle X^2 \rangle - \langle X \rangle^2,
	\qquad
	\langle X \rangle = \operatorname{Tr}[\hat{\rho} X].
\end{equation}
For any pair of self-adjoint operators $X$ and $Y$ the Robertson relation holds,
\begin{equation}
	\Delta X \, \Delta Y \geq \frac{1}{2} \left| \langle [X,Y] \rangle \right|,
\end{equation}
which reduces to the familiar Heisenberg uncertainty relation for position and momentum. 
Uncertainty relations express fundamental limitations on the simultaneous sharpness of non-commuting observables, not a limitation of measurement technology.

\subsection{Generalized measurements and POVMs}
\label{sec:povm}

Ideal projective measurements are an important but special case. 
Realistic measurements often involve finite resolution, inefficiencies or measurements on a subsystem after coupling to an ancilla. 
The general structure is captured by positive operator-valued measures (POVMs).

A POVM is a set of positive operators $\{E_x\}$ on $\mathcal{H}$ satisfying
\begin{equation}
	E_x \geq 0,
	\qquad
	\sum_x E_x = \mathbb{1}.
\end{equation}
The probability of observing outcome $x$ when the system is in the state $\hat{\rho}$ is
\begin{equation}
	p(x) = \operatorname{Tr}[\hat{\rho} E_x].
\end{equation}
Projective measurements correspond to the special case in which the $E_x$ are orthogonal projectors. 
In general, POVM elements need not be projectors, nor mutually orthogonal, and can have rank larger than one.

POVMs arise naturally from the \emph{indirect measurement} scheme, in which the system interacts with an ancillary probe, and a projective measurement is performed only on the probe. 
By tracing out the probe degrees of freedom, the effective action on the system is described by a POVM.

\subsection{Measurement operators and state update}

To account for both statistics and state update, it is convenient to introduce \emph{measurement operators} (also called Kraus operators) $\{M_x\}$ acting on $\mathcal{H}$, such that
\begin{equation}
	E_x = M_x^\dagger M_x,
	\qquad
	\sum_x M_x^\dagger M_x = \mathbb{1}.
\end{equation}
The probability of outcome $x$ is again $p(x) = \operatorname{Tr}[\hat{\rho} E_x]$, and the post-measurement state conditioned on $x$ is
\begin{equation}
	\hat{\rho} \longrightarrow \hat{\rho}_x =
	\frac{M_x \hat{\rho} M_x^\dagger}{p(x)}.
\end{equation}
Different sets of operators $\{M_x\}$ may produce the same POVM $\{E_x\}$ but correspond to different physical measurement schemes (for example, different couplings to an ancilla or different post-processing of the measurement outcomes).

\subsection{Naimark dilation and indirect measurements}

A fundamental structural result is the Naimark dilation theorem: any POVM on $\mathcal{H}_S$ can be realized as a projective measurement on a larger Hilbert space.

Formally, given a POVM $\{E_x\}$ on $\mathcal{H}_S$, there exist an ancillary Hilbert space $\mathcal{H}_E$, an initial ancilla state $\ket{0}_E \in \mathcal{H}_E$, a unitary $U$ on $\mathcal{H}_S \otimes \mathcal{H}_E$ and a family of projectors $\{P_x\}$ on $\mathcal{H}_E$ such that
\begin{equation}
	E_x =
	\operatorname{Tr}_E \left[
	U^\dagger (\mathbb{1}_S \otimes P_x) U
	(\mathbb{1}_S \otimes \ket{0}_E\bra{0})
	\right].
\end{equation}
Physical detectors are thus naturally modelled as quantum systems interacting with the system of interest; the POVM formalism arises when one describes only the effect of this interaction on the system.

\section{Quantum dynamics and quantum channels}
\label{sec:channels}

\subsection{Unitary dynamics of closed systems}

For an isolated system with Hamiltonian $H$, the state vector obeys the Schrödinger equation,
\begin{equation}
	i \hbar \frac{\mathrm{d}}{\mathrm{d}t} \ket{\psi(t)} = H \ket{\psi(t)}.
\end{equation}
For a time-independent Hamiltonian, the formal solution is
\begin{equation}
	\ket{\psi(t)} = U(t,t_0)\ket{\psi(t_0)},
	\qquad
	U(t,t_0) = \exp\left[-\frac{i}{\hbar} H (t-t_0)\right],
\end{equation}
with $U$ unitary, $U^\dagger U = UU^\dagger = \mathbb{1}$.

In the density matrix formalism, the dynamics is governed by the Liouville--von Neumann equation,
\begin{equation}
	i \hbar \frac{\mathrm{d}\hat{\rho}(t)}{\mathrm{d}t} = [H,\hat{\rho}(t)],
\end{equation}
whose solution is
\begin{equation}
	\hat{\rho}(t) = U(t,t_0) \hat{\rho}(t_0) U^\dagger(t,t_0).
\end{equation}
For time-dependent Hamiltonians, the solution involves the time-ordered exponential, but the evolution map remains unitary.

\subsection{Schrödinger, Heisenberg and interaction pictures}

There are three frequently used pictures of quantum dynamics:

\begin{itemize}
	\item In the \emph{Schrödinger picture}, states carry the time dependence and operators are fixed (unless they have explicit time dependence).
	\item In the \emph{Heisenberg picture}, states are fixed and operators evolve according to
	\begin{equation}
		X_H(t) = U^\dagger(t,t_0) X_S U(t,t_0),
	\end{equation}
	with equation of motion
	\begin{equation}
		i \hbar \frac{\mathrm{d}}{\mathrm{d}t} X_H(t) = [X_H(t), H_H(t)] + i \hbar \left(\frac{\partial X}{\partial t}\right)_H.
	\end{equation}
	\item In the \emph{interaction picture}, the Hamiltonian is split as $H = H_0 + H_I(t)$. 
	The evolution due to $H_0$ is treated exactly, while $H_I$ is considered as an interaction. 
	This picture is particularly convenient for time-dependent perturbation theory and for deriving effective dynamics in open systems.
\end{itemize}
Despite their different appearance, all pictures are physically equivalent, and the choice is usually dictated by convenience.

\subsection{Open systems and completely positive maps}

Real systems are never perfectly isolated: they interact with external degrees of freedom, typically modelled as an environment or bath. 
Let us denote by $S$ the system of interest and by $E$ the environment, with Hilbert spaces $\mathcal{H}_S$ and $\mathcal{H}_E$. 
The total Hilbert space is $\mathcal{H}_S \otimes \mathcal{H}_E$, and the joint state $\hat{\rho}_{SE}$ evolves unitarily under a total Hamiltonian $H_{SE}$, giving rise to a unitary propagator $U_{SE}$.

If at time $t_0$ the joint state factorizes as
\begin{equation}
	\hat{\rho}_{SE}(t_0) = \hat{\rho}_S(t_0) \otimes \hat{\rho}_E,
\end{equation}
then at a later time $t$ the reduced state of the system is
\begin{equation}
	\hat{\rho}_S(t) 
	= 
	\operatorname{Tr}_E\left[
	U_{SE}(t,t_0) \bigl( \hat{\rho}_S(t_0) \otimes \hat{\rho}_E \bigr) U_{SE}^\dagger(t,t_0)
	\right].
\end{equation}
This defines a map $\mathcal{E}$ from initial to final system states,
\begin{equation}
	\hat{\rho}_S(t) = \mathcal{E}\bigl(\hat{\rho}_S(t_0)\bigr).
\end{equation}
Physically admissible maps of this sort must preserve positivity and trace, and must remain positive when extended to larger Hilbert spaces. 
This leads to the notion of completely positive trace-preserving (CPTP) maps.

\subsection{Completely positive trace-preserving maps and Kraus representation}

A linear map $\mathcal{E} : \mathcal{L}(\mathcal{H}_S) \to \mathcal{L}(\mathcal{H}_S)$ is
\begin{itemize}
	\item \emph{trace-preserving} if $\operatorname{Tr}[\mathcal{E}(\hat{\rho})] = \operatorname{Tr}[\hat{\rho}]$ for all $\hat{\rho}$;
	\item \emph{positive} if $\hat{\rho} \geq 0 \Rightarrow \mathcal{E}(\hat{\rho}) \geq 0$;
	\item \emph{completely positive} if, for any ancillary Hilbert space $\mathcal{H}_A$ and any positive operator $\hat{\sigma}$ on $\mathcal{H}_S \otimes \mathcal{H}_A$, we have
	\begin{equation}
		(\mathcal{E} \otimes \mathbb{1}_A)(\hat{\sigma}) \geq 0.
	\end{equation}
\end{itemize}
The last requirement guarantees that the map is physically meaningful even when the system is entangled with some ancilla.

A fundamental theorem (Stinespring--Kraus--Choi--Sudarshan) states that a map $\mathcal{E}$ is CPTP if and only if it admits a \emph{Kraus representation}:
\begin{equation}
	\mathcal{E}(\hat{\rho}) = \sum_k K_k \hat{\rho} K_k^\dagger,
	\qquad
	\sum_k K_k^\dagger K_k = \mathbb{1}_S,
\end{equation}
where the $K_k$ are operators on $\mathcal{H}_S$, called Kraus operators. 
The trace-preserving condition is encoded in the normalization $\sum_k K_k^\dagger K_k = \mathbb{1}_S$.

If the system initially interacts unitarily with an environment and then the environment is traced out, the resulting map on the system is always CPTP. Conversely, any CPTP map can be realized in this way for a suitable choice of environment, initial environmental state and joint unitary.

\subsection{Examples: dephasing and amplitude damping}

\paragraph{Dephasing channel.}
Consider a single qubit with computational basis $\{\ket{0}, \ket{1}\}$.
A pure dephasing channel with dephasing probability $p$ can be written as
\begin{equation}
	\mathcal{E}_{\text{deph}}(\hat{\rho})
	=
	(1-p) \hat{\rho} + p \, \sigma_z \hat{\rho} \sigma_z,
\end{equation}
where $\sigma_z$ is the Pauli operator with eigenstates $\ket{1}$ (eigenvalue $+1$) and $\ket{0}$ (eigenvalue $-1$).
A Kraus representation is given by
\begin{equation}
	K_0 = \sqrt{1-p}\, \mathbb{1}_2,
	\qquad
	K_1 = \sqrt{p}\, \sigma_z.
\end{equation}
This channel leaves populations invariant while suppressing the off-diagonal coherences in the computational basis.

\paragraph{Amplitude damping channel.}
The amplitude damping channel describes spontaneous decay from the excited state $\ket{1}$ to the ground state $\ket{0}$ with probability $\gamma$. 
The Kraus operators can be chosen as
\begin{equation}
	K_0 =
	\begin{pmatrix}
		1 & 0 \\
		0 & \sqrt{1-\gamma}
	\end{pmatrix},
	\qquad
	K_1 =
	\begin{pmatrix}
		0 & \sqrt{\gamma} \\
		0 & 0
	\end{pmatrix}
\end{equation}
in the basis $\{\ket{0},\ket{1}\}$. 
This model will reappear in the description of dissipative processes in open quantum systems and in the construction of master equations.

\section{Qubits, spin-\texorpdfstring{$1/2$}{1/2} systems and Bloch sphere}

\subsection{Qubits and Pauli operators}

The simplest quantum system is the two-level system or \emph{qubit}. 
In this thesis we adopt the convention in which $\ket{0}$ denotes the ground state and $\ket{1}$ the excited state, with
\begin{equation}
	\sigma_z \ket{1} = + \ket{1},
	\qquad
	\sigma_z \ket{0} = - \ket{0}.
\end{equation}
In the ordered basis $\{\ket{1},\ket{0}\}$ the Pauli operators read
\begin{equation}
	\sigma_x = 
	\begin{pmatrix}
		0 & 1 \\
		1 & 0
	\end{pmatrix},
	\quad
	\sigma_y =
	\begin{pmatrix}
		0 & -i \\
		i & 0
	\end{pmatrix},
	\quad
	\sigma_z =
	\begin{pmatrix}
		1 & 0 \\
		0 & -1
	\end{pmatrix},
\end{equation}
and the identity is
\begin{equation}
	\mathbb{1}_2 =
	\begin{pmatrix}
		1 & 0 \\
		0 & 1
	\end{pmatrix}.
\end{equation}
The eigenstates of $\sigma_x$ are
\begin{equation}
	\ket{+} = \frac{1}{\sqrt{2}}\bigl(\ket{1} + \ket{0}\bigr),
	\qquad
	\ket{-} = \frac{1}{\sqrt{2}}\bigl(\ket{1} - \ket{0}\bigr),
\end{equation}
and analogous expressions hold for the eigenstates of $\sigma_y$.

A spin-$1/2$ system is physically equivalent to a qubit: its spin operators are related to the Pauli operators by
\begin{equation}
	S_\alpha = \frac{\hbar}{2} \sigma_\alpha,
	\qquad
	\alpha = x,y,z.
\end{equation}

\subsection{Bloch sphere representation}

Any single-qubit density operator can be written as
\begin{equation}
	\hat{\rho} = \frac{1}{2}\left(\mathbb{1}_2 + \mathbf{r} \cdot \boldsymbol{\sigma}\right),
\end{equation}
where $\mathbf{r} = (r_x,r_y,r_z)$ is a real vector with $|\mathbf{r}| \leq 1$, and $\boldsymbol{\sigma} = (\sigma_x,\sigma_y,\sigma_z)$. 
The vector $\mathbf{r}$ is called the \emph{Bloch vector}. 
Pure states correspond to $|\mathbf{r}| = 1$ and lie on the surface of the Bloch sphere; mixed states correspond to points inside the sphere, with the maximally mixed state at the origin.

In spherical coordinates, a pure state can be written as
\begin{equation}
	\ket{\psi(\theta,\phi)} = \cos\frac{\theta}{2} \ket{0}
	+ e^{i\phi} \sin\frac{\theta}{2} \ket{1},
\end{equation}
which corresponds to a Bloch vector
\begin{equation}
	\mathbf{r} = 
	(\sin\theta \cos\phi, \, \sin\theta \sin\phi, \, \cos\theta).
\end{equation}
Unitary dynamics generated by a Hamiltonian proportional to a Pauli operator,
\begin{equation}
	H = \frac{\hbar \Omega}{2} \, \mathbf{n} \cdot \boldsymbol{\sigma},
\end{equation}
corresponds to rotations of the Bloch vector around the axis $\mathbf{n}$ with angular frequency $\Omega$.

\subsection{Multi-qubit systems, Bell states and entanglement}

For two qubits, the Hilbert space is $\mathbb{C}^2 \otimes \mathbb{C}^2$. 
A convenient basis is the computational basis
\begin{equation}
	\left\{\ket{11}, \ket{10}, \ket{01}, \ket{00} \right\},
\end{equation}
where $\ket{ab} = \ket{a}_1 \otimes \ket{b}_2$ and we keep the convention that $\ket{1}$ is the excited state and $\ket{0}$ the ground state.

The four Bell states form an orthonormal basis of maximally entangled states,
\begin{align}
	\ket{\phi^\pm} &= \frac{1}{\sqrt{2}}\left(\ket{00} \pm \ket{11}\right),
	\\
	\ket{\psi^\pm} &= \frac{1}{\sqrt{2}}\left(\ket{01} \pm \ket{10}\right).
\end{align}
A pure bipartite state $\ket{\psi_{AB}}$ is called \emph{separable} if it can be written as a product $\ket{\psi_A} \otimes \ket{\psi_B}$. 
Otherwise it is entangled. 
The Bell states are canonical examples of entangled states.

For mixed states, separability is defined by the requirement that the density operator can be written as a convex combination of product states,
\begin{equation}
	\hat{\rho}_{AB} = \sum_k p_k \, \hat{\rho}_A^{(k)} \otimes \hat{\rho}_B^{(k)}.
\end{equation}
If such a decomposition does not exist, the state is entangled. 
Quantifying entanglement requires suitable measures, such as concurrence for two-qubit systems, which will be employed in later chapters.

\section{Next chapter}

In this chapter we have introduced the mathematical framework of quantum mechanics in the language used in the rest of the thesis. 
These concepts provide the foundation for the description of open quantum systems and continuously monitored dynamics. 
In the next chapter we will build on this formalism to introduce master equations, the Markovian approximation and, subsequently, stochastic master equations describing the conditional evolution of systems under continuous measurements.
